\documentclass[10pt]{article}

% ---------- Document Identity (update these only) ----------
\newcommand{\DocStage}{}
\newcommand{\DocTitle}{Your Road to Tier One \textbf{SOF}}
\newcommand{\ProgramName}{182-Week Civilian-to-Selection Readiness Pipeline}
\newcommand{\ProgramSubtitle}{}

% ---------- Page ----------
\usepackage[legalpaper,left=0.5in,right=0.5in,top=0.6in,bottom=0.45in]{geometry}

% ---------- Typography ----------
\usepackage[T1]{fontenc}
\usepackage[utf8]{inputenc}
\usepackage{libertinus}
\usepackage{microtype}
\usepackage{parskip}
\usepackage{ragged2e}
\usepackage{textcomp}
\AtBeginDocument{\color{black}}
\AtBeginDocument{\pagecolor{white}}
\AtBeginDocument{\justifying}

% ---------- Landscape pages ----------
\usepackage{pdflscape}

% ---------- Color ----------
\usepackage[table]{xcolor}
\definecolor{StageBlue}{HTML}{1F4E79}
\definecolor{StageBlueDark}{HTML}{163A5A}
\definecolor{LightGray}{HTML}{F2F4F7}
\definecolor{MidGray}{HTML}{D9DEE5}
\definecolor{RuleGray}{HTML}{C7CED8}
\definecolor{HeaderInk}{HTML}{000000}
\definecolor{Ink}{HTML}{000000}

% ---------- Utility ----------
\usepackage{enumitem}
\sloppy
\setlength{\emergencystretch}{2em}

% ---------- Boxes / UI ----------
\usepackage[most]{tcolorbox}

\tcbset{
  charterTitle/.style={
    enhanced,
    colback=StageBlue!4,
    colframe=StageBlue,
    boxrule=1.25pt,
    arc=3pt,
    left=16pt,right=16pt,top=14pt,bottom=14pt,
    borderline north={3.2pt}{0pt}{StageBlue},
    borderline west={4pt}{0pt}{StageBlue},
    borderline south={1.25pt}{0pt}{StageBlue},
    drop shadow={black!25!white},
  },
  charterRule/.style={
    enhanced,
    colback=LightGray,
    colframe=RuleGray,
    boxrule=0.85pt,
    arc=3pt,
    left=12pt,right=12pt,top=10pt,bottom=10pt,
    borderline west={3pt}{0pt}{StageBlue},
  },
  charterBadge/.style={
    enhanced,
    colback=StageBlue,
    coltext=white,
    colframe=StageBlueDark,
    boxrule=0.6pt,
    arc=2pt,
    left=6pt,right=6pt,top=3pt,bottom=3pt,
  }
}
\newtcbox{\Badge}{nobeforeafter,charterBadge}

% ---------- Lists ----------
\setlist[itemize]{leftmargin=1.5em, itemsep=0.25em, topsep=0.25em}
\setlist[enumerate]{leftmargin=1.8em, itemsep=0.25em, topsep=0.25em}

% ---------- TOC ----------
\usepackage{tocloft}
\setcounter{tocdepth}{3}
\setlength{\cftbeforesecskip}{3pt}
\setlength{\cftbeforesubsecskip}{2pt}
\setlength{\cftbeforesubsubsecskip}{0pt}
\renewcommand{\cftsecfont}{\bfseries}
\renewcommand{\cftsecpagefont}{\bfseries}
\renewcommand{\cftsubsecfont}{\normalfont}
\renewcommand{\cftsubsecpagefont}{\normalfont}
\renewcommand{\cftsubsubsecfont}{\normalfont}
\renewcommand{\cftsubsubsecpagefont}{\normalfont}
\setlength{\cftsecindent}{0pt}
\setlength{\cftsubsecindent}{1.25em}
\setlength{\cftsubsubsecindent}{2.8em}
\setlength{\cftsecnumwidth}{2.1em}
\setlength{\cftsubsecnumwidth}{2.8em}
\setlength{\cftsubsubsecnumwidth}{3.6em}

% Remove the built-in TOC heading so only the custom heading is shown
\makeatletter
\renewcommand{\tableofcontents}{\@starttoc{toc}}
\makeatother

% ---------- Header/Footer: page number only ----------
\usepackage{fancyhdr}
\pagestyle{fancy}
\fancyhf{}
\renewcommand{\headrulewidth}{0pt}
\renewcommand{\footrulewidth}{0pt}
\fancyfoot[C]{\thepage}

% ---------- Section formatting ----------
\usepackage{titlesec}

\titleformat{\section}
  {\Large\bfseries\color{Ink}}
  {\thesection}{0.6em}{}
  [\vspace{0.25em}\textcolor{StageBlue}{\rule{\linewidth}{0.8pt}}]

\titleformat{\subsection}
  {\large\bfseries\color{Ink}}
  {\thesubsection}{0.6em}{}

\titleformat{\subsubsection}
  {\normalsize\bfseries\color{Ink}}
  {\thesubsubsection}{0.6em}{}

\titlespacing*{\section}{0pt}{0.95em}{0.35em}
\titlespacing*{\subsection}{0pt}{0.65em}{0.25em}
\titlespacing*{\subsubsection}{0pt}{0.55em}{0.2em}

% ---------- Table engine ----------
\usepackage{tabularray}
\UseTblrLibrary{booktabs}

% ---------- tabularray: GLOBAL table treatment ----------
\SetTblrInner{
  cells = {fg=black, bg=white, valign=t},
}

% ---------- Header helpers ----------
\newcommand{\Hdr}[1]{\shortstack[c]{\strut #1\strut}}

% ---------- Lines ----------
\newcommand{\TitleLine}{\textcolor{StageBlue}{\rule{0.98\linewidth}{0.95pt}}}
\newcommand{\SubTitleLine}{\textcolor{RuleGray}{\rule{0.98\linewidth}{0.65pt}}}

% ---------- Continued on next page ----------
\newcommand{\ContinuedOnNextPageInline}{%
  \par
  \vfill
  \noindent\textcolor{RuleGray}{\rule{\linewidth}{1pt}}\vspace{-2pt}
  \noindent\hfill\textit{Continued on next page}\par\vspace{-10pt}
  \noindent\textcolor{RuleGray}{\rule{\linewidth}{1pt}}
  \clearpage
}

% ---------- PDF hyperlinks ----------
\usepackage[hidelinks]{hyperref}
\hypersetup{
  colorlinks=false,
  pdfborder={0 0 0},
  linktoc=all,
  pdftitle={\DocTitle},
  pdfauthor={\ProgramName}
}

% =========================================================
\begin{document}

\section*{Stage 11.C — Recovery and Injury Prevention in Final Phases}
\phantomsection
\addcontentsline{toc}{section}{Stage 11.C — Recovery and Injury Prevention in Final Phases}

\subsection*{11.C.1 Purpose and Scope}
\phantomsection
\addcontentsline{toc}{subsection}{11.C.1 Purpose and Scope}

11.C exists to preserve readiness in the final phases by reducing preventable injury risk, limiting validity drift, and sustaining repeatable performance within protected Deload/Test windows. In the late pipeline, consequence is high: small breakdowns (sleep collapse, hotspots progressing to blisters, gait drift, illness) can invalidate tests, force reslots, and destabilize gate decisions.

11.C does not provide medical treatment guidance, rehabilitation protocols, or programming detail. It does not prescribe workouts, sets/reps, interval prescriptions, weekly splits, or day-by-day schedules. It provides governance overlays, risk surveillance, and operational controls aligned to Stage 1.F validity doctrine, Stage 3.E adjustment posture, and Stage 2.E logging discipline.

\newpage

\subsection*{11.C.2 Recovery Governance Hierarchy}
\phantomsection
\addcontentsline{toc}{subsection}{11.C.2 Recovery Governance Hierarchy}

Priority order (binding):

\begin{itemize}
\item Safety and stop posture override schedule and performance:
\begin{itemize}
\item \textbf{STOP} \textbf{NOW} and \textbf{MEDICAL} \textbf{HOLD} posture applies when indicated; do not negotiate.
\end{itemize}
\item Validity posture overrides “good effort”:
\begin{itemize}
\item inadmissible evidence cannot support \textbf{ADVANCE}.
\end{itemize}
\item Recovery actions are tools to preserve admissible evidence and reduce risk:
\begin{itemize}
\item recovery actions are not a license to force progression or to compensate for missing compliance by stacking.
\end{itemize}
\end{itemize}

Do not do list (binding):

\begin{itemize}
\item Do not train through gait-altering pain.
\item Do not stack make-up tests or compress missed tests into a single week.
\item Do not add novelty near gates or Deload/Test weeks (new shoes, new route class, new tools, new hydration/electrolyte practices, major nutrition structure changes).
\end{itemize}

\newpage

\subsection*{11.C.3 Final-Phase Risk Profile and What Changes Late Pipeline}
\phantomsection
\addcontentsline{toc}{subsection}{11.C.3 Final-Phase Risk Profile and What Changes Late Pipeline}

Late phases increase consequence because:

\begin{itemize}
\item specificity and consequence of gate-relevant tests increase,
\item protected windows concentrate decision-grade evidence demands,
\item fatigue tolerance is narrower and error costs are higher,
\item environment volatility and access constraints (route/pool/tool availability) can invalidate otherwise strong performance.
\end{itemize}

Key risk families to control (Stage 1.B references):

\begin{itemize}
\item Impact and overuse drift in \textbf{RUN} exposure — \textbf{RISK-1B-001}.
\item Load-bearing tolerance failures and interface breakdown (ruck fit, load shift) — \textbf{RISK-1B-001}.
\item Footwear/footcare failures leading to hotspots → blisters → gait drift — \textbf{RISK-1B-009}.
\item Sleep debt and recovery collapse contaminating validity and increasing injury risk — \textbf{RISK-1B-004}.
\item Environment hazards (ice/heat/air quality/visibility) driving unsafe or invalid test conditions — \textbf{RISK-1B-008} / \textbf{RISK-1B-013}.
\item Test invalidity due to missing required fields, tool drift, or environment \textbf{ID} gaps — \textbf{RISK-1B-014}.
\item Water governance breaches (any unsupervised underwater/breath-hold attempt) — \textbf{RISK-1B-010}.
\end{itemize}

\newpage
    
\subsection*{11.C.4 Recovery Surveillance Signals and Triggers}
\phantomsection
\addcontentsline{toc}{subsection}{11.C.4 Recovery Surveillance Signals and Triggers}

Decision signals are drawn from existing Stage 2 metrics/logging; 11.C does not create a new metric system. Signals drive downshifts, reslots, and conservative gate posture.

Signal families:

\begin{itemize}
\item \textbf{RECOVERY} signals:
\begin{itemize}
\item sleep hours and sleep quality trends,
\item fatigue and soreness trends.
\end{itemize}
\item \textbf{DURABILITY} signals:
\begin{itemize}
\item pain severity and location,
\item gait-altering pain flag,
\item foot hotspot and blister flags.
\end{itemize}
\item \textbf{COMPLIANCE} and validity signals:
\begin{itemize}
\item repeated \textbf{MODIFIED} / \textbf{INVALID} sessions,
\item missed logs or missing required fields (\textbf{UNKNOWN} / \textbf{MISSING} recorded; no inference).
\end{itemize}
\item Illness flags:
\begin{itemize}
\item non-diagnostic illness posture triggers; do not test through illness.
\end{itemize}
\end{itemize}

\begingroup
\scriptsize
\setlength{\tabcolsep}{2pt}
\renewcommand{\arraystretch}{1.12}

\begin{longtblr}{
width=\linewidth,
rowhead=1,
colspec = {
X[l,1.35]
X[l,1.55]
X[l,1.55]
X[l,1.15]
X[l,1.15]
},
hlines,
vlines,
cells = {hsep=6pt, vsep=6pt, halign=l, valign=t},
row{odd} = {bg=white},
row{even} = {bg=LightGray},
row{1} = {bg=StageBlue!15, fg=black, font=\bfseries, halign=l, valign=m},
}
\Hdr{Signal Cluster} &
\Hdr{Examples} &
\Hdr{Default Action} &
\Hdr{Validity Impact} &
\Hdr{What to Log} \\
Sleep debt / recovery strain &
sleep hours trending down; sleep quality trending low; fatigue rising &
Downshift exposure intensity/density; consider maintenance reset posture via Stage 10 rules; protect protected windows &
If severe near gates: convert to check-in-only or reslot; do not chase maximal attempts &
Sleep markers; fatigue/soreness; decision; hold period; gate proximity \\
Pain and mechanics drift &
pain increasing; pain alters mechanics; gait-altering pain flag Yes &
\textbf{STOP} \textbf{NOW} for that exposure; substitute lower-risk modality; consider \textbf{MEDICAL} \textbf{HOLD} if indicated &
Gate attempts requiring \textbf{RUN}/\textbf{RUCK} are disqualified until mechanics stable; reslot tests &
Pain severity/location; gait flag; foot flags; exposure stopped/substituted \\
Foot integrity breakdown &
hotspot progressing; blister present; repeated hotspot flags &
Downshift friction exposure; enforce foot system stabilization; avoid long \textbf{RUN}/\textbf{RUCK} attempts &
If blister/gait drift present: reslot gate tests; do not claim comparable evidence &
Hotspot/blister flag + location; footwear/sock notes; substitution logged \\
Validity drift / admin failure &
missing required fields; late entries; repeated \textbf{INVALID} sessions &
Restore logging discipline; simplify to minimum viable session posture without skipping required elements &
Missing fields → inadmissible evidence; default \textbf{HOLD} if gate evidence incomplete &
Missing fields recorded as \textbf{MISSING}; validity tags; reason posture; corrective action \\
Illness flags &
fever/systemic symptoms; respiratory illness; severe GI distress &
Suspend maximal efforts; reslot tests; \textbf{MEDICAL} \textbf{HOLD} when indicated &
Testing during illness is inadmissible; reslot per Stage 3 &
Illness flags; symptoms notes; reslot decision; \textbf{MEDICAL} \textbf{HOLD} if applied \\
\end{longtblr}

\endgroup

\newpage

\subsection*{11.C.5 Prevention Doctrine by Exposure Type}
\phantomsection
\addcontentsline{toc}{subsection}{11.C.5 Prevention Doctrine by Exposure Type}

This section defines operational controls, stop triggers, downshift posture, and documentation requirements. It does not prescribe workouts.

\subsubsection*{\textbf{RUN} exposure controls}
\phantomsection

\begin{itemize}
\item Impact-managed posture:
\begin{itemize}
\item \textbf{RUN} escalation \textbf{MUST} be denied when gait-altering pain is present or when foot flags indicate hotspot→blister progression risk.
\end{itemize}
\item Footwear interface must be stable and logged:
\begin{itemize}
\item footwear system changes near protected windows are prohibited; document any unavoidable change and treat outcomes conservatively.
\end{itemize}
\item Environment hazard checks:
\begin{itemize}
\item ice glaze, heat index concern, air quality warnings, and poor visibility trigger indoor substitution or reslot.
\end{itemize}
\end{itemize}

Stop triggers:

\begin{itemize}
\item gait-altering pain, chest pain/pressure, syncope/near-syncope, disproportionate dyspnea, new neurologic symptoms, heat illness signs → \textbf{STOP} \textbf{NOW}; \textbf{MEDICAL} \textbf{HOLD} when indicated.
\end{itemize}

Downshift posture:

\begin{itemize}
\item substitute to lower-impact conditioning modality when \textbf{RUN} is unsafe; preserve cadence where feasible.
\end{itemize}

Documentation:

\begin{itemize}
\item Route/Machine identity and conditions; footwear notes; validity tag and reason posture if modified.
\end{itemize}

\subsubsection*{\textbf{RUCK} and load-bearing exposure controls}
\phantomsection

\begin{itemize}
\item Load verification posture and pack fit checks:
\begin{itemize}
\item dry load recorded; water logged separately; ruck fit stable; load shift is prohibited.
\end{itemize}
\item No novelty pack/interface changes near gates:
\begin{itemize}
\item new ruck interface or new load module changes are prohibited in protected windows.
\end{itemize}
\end{itemize}

Stop triggers:

\begin{itemize}
\item gait-altering pain; hotspot progression; instability; red-flag symptoms → \textbf{STOP} \textbf{NOW}.
\end{itemize}

Downshift posture:

\begin{itemize}
\item reduce risk exposure; substitute with lower-impact modality; reslot ruck tests; do not “prove it.”
\end{itemize}

Documentation:

\begin{itemize}
\item Route \textbf{ID}; dry load; carried water; foot flags; visibility conditions; validity tag and reason posture.
\end{itemize}

\subsubsection*{\textbf{CAL} volume controls}
\phantomsection

\begin{itemize}
\item Tissue tolerance posture:
\begin{itemize}
\item avoid sudden volume spikes; avoid novelty equipment changes (bars, surfaces) near protected windows.
\end{itemize}
\item Grip/elbow/shoulder caution posture:
\begin{itemize}
\item discomfort that changes mechanics triggers downshift; do not chase “max reps” outside protected windows.
\end{itemize}
\end{itemize}

Stop triggers:

\begin{itemize}
\item acute pain that alters mechanics; neurologic symptoms; red flags → \textbf{STOP} \textbf{NOW}.
\end{itemize}

Downshift posture:

\begin{itemize}
\item preserve cadence with reduced risk exposure; use conservative substitutions; do not force maximal attempts.
\end{itemize}

Documentation:

\begin{itemize}
\item tool/surface notes; evidence artifact posture where required; validity tag.
\end{itemize}

\subsubsection*{\textbf{STR} exposure controls}
\phantomsection

\begin{itemize}
\item Submax posture near gates:
\begin{itemize}
\item avoid novelty PR chasing and unsafe attempts in proximity to high-consequence test windows.
\end{itemize}
\item Implement identity stability:
\begin{itemize}
\item if implement changes, log it and interpret trends conservatively.
\end{itemize}
\end{itemize}

Stop triggers:

\begin{itemize}
\item acute pain/instability; red-flag symptoms → \textbf{STOP} \textbf{NOW}.
\end{itemize}

Downshift posture:

\begin{itemize}
\item reduce intensity/density; preserve cadence with safer exposures.
\end{itemize}

Documentation:

\begin{itemize}
\item tool/implement identity and load method; validity tag.
\end{itemize}

\subsubsection*{\textbf{SWIM} exposure controls}
\phantomsection

\begin{itemize}
\item Surface-only unless governed:
\begin{itemize}
\item underwater/breath-hold exposure is prohibited unless governance prerequisites are met; otherwise not scheduled.
\end{itemize}
\item Pool conditions must support validity:
\begin{itemize}
\item Pool \textbf{ID} and length unit logged; unsafe/crowded conditions that compromise validity trigger reslot.
\end{itemize}
\end{itemize}

Stop triggers:

\begin{itemize}
\item distress, disorientation, panic, dizziness, cramps, illness symptoms → \textbf{STOP} \textbf{NOW}; exit water; seek evaluation posture when indicated.
\end{itemize}

Downshift posture:

\begin{itemize}
\item reslot swim tests; substitute to non-aquatic conditioning without claiming equivalence for gate evidence.
\end{itemize}

Documentation:

\begin{itemize}
\item Pool \textbf{ID} + length unit; conditions notes; validity tag; Evidence Ref posture as required.
\end{itemize}

\subsubsection*{Environment controls (all exposures)}
\phantomsection

\begin{itemize}
\item Heat/ice/air quality/visibility substitution posture:
\begin{itemize}
\item unsafe conditions trigger substitution indoors or reslot per Stage 3; do not force.
\end{itemize}
\end{itemize}

Stop triggers:

\begin{itemize}
\item lightning, ice glaze, whiteout, official air quality warning, heat illness signs → \textbf{STOP} \textbf{NOW} for outdoor exposure.
\end{itemize}

Documentation:

\begin{itemize}
\item environment notes; substitution/reslot reason; validity impact.
\end{itemize}

\newpage

\subsection*{11.C.6 Downshift and Re-Entry Posture}
\phantomsection
\addcontentsline{toc}{subsection}{11.C.6 Downshift and Re-Entry Posture}

Allowed actions consistent with Stage 3.E (binding):

\begin{itemize}
\item Preserve cadence where feasible:
\begin{itemize}
\item downshift exposure risk and intensity; do not drop the schedule as the first lever.
\end{itemize}
\item Substitute to lower-risk modalities:
\begin{itemize}
\item substitutions must preserve intent and must be logged.
\end{itemize}
\item Reslot tests to next protected window:
\begin{itemize}
\item invalid or unsafe tests are reslotted; no stacking make-ups.
\end{itemize}
\item Default gate outcomes:
\begin{itemize}
\item \textbf{HOLD} when gate evidence is incomplete or inadmissible,
\item \textbf{MEDICAL} \textbf{HOLD} when red flags or unsafe symptoms appear per Stage 1.F posture.
\end{itemize}
\end{itemize}

Decision table:

\begingroup
\scriptsize
\setlength{\tabcolsep}{2pt}
\renewcommand{\arraystretch}{1.12}

\begin{longtblr}{
width=\linewidth,
colspec = {
X[c,1.65]
X[c,2.25]
X[c,2.85]
X[c,2.85]
X[c,1.60]
},
rowhead=1,
hlines,
vlines,
cells = {hsep=2pt, vsep=6pt, halign=c, valign=m},
row{1} = {bg=StageBlue!15, fg=black, font=\bfseries, valign=m},
row{odd} = {bg=white},
row{even} = {bg=LightGray},
}
\Hdr{Situation} &
\Hdr{Example Trigger} &
\Hdr{Immediate Action} &
\Hdr{Next-Step Adjustment} &
\Hdr{Gate Posture} \\
Gait-altering pain appears &
limp/compensation noted; gait flag Yes &
\textbf{STOP} \textbf{NOW} for \textbf{RUN}/\textbf{RUCK} exposure; substitute lower-risk modality &
Reslot any \textbf{RUN}/\textbf{RUCK} tests; maintain cadence with safer exposures &
\textbf{HOLD} (default) or \textbf{MEDICAL} \textbf{HOLD} if indicated \\
Hotspot trending toward blister &
hotspot flag repeats; pain increasing &
Downshift friction exposure; stabilize footwear/sock interface &
Reslot high-cost tests; enforce foot protocol consistency &
\textbf{HOLD} (default) \\
Sleep debt + fatigue trend &
sleep hours/quality down; fatigue high &
Downshift intensity/density; hold deficit posture per Stage 10; consider maintenance reset per Stage 10 rules &
Convert planned test week to deload-only/check-in-only if needed; reslot maximal tests &
\textbf{HOLD} (default) \\
Illness flags present &
fever/systemic symptoms; severe GI distress &
\textbf{STOP} \textbf{NOW} for maximal attempts; consider \textbf{MEDICAL} \textbf{HOLD} &
Reslot tests; preserve safety posture; no testing through illness &
\textbf{MEDICAL} \textbf{HOLD} when indicated \\
Environment invalidates safety &
ice glaze, lightning, air quality warning &
Cancel outdoor attempt immediately &
Substitute indoors only if comparable; otherwise reslot &
\textbf{HOLD} (default) \\
Logging completeness failure &
missing required fields; late entry &
Treat evidence as inadmissible; correct logging system &
Restore same-day capture; do not infer missing data &
\textbf{HOLD} (default) \\
\end{longtblr}

\endgroup

Requires Stage 8 review flags (if encountered):

\begin{itemize}
\item Any adjustment approach that attempts to stack make-ups to "catch up" must be flagged.
\item Any attempt to treat inadmissible evidence (missing required fields) as gate-grade evidence.
\end{itemize}

\newpage

\subsection*{11.C.7 Tooling and Environment Supports for Recovery Execution}
\phantomsection
\addcontentsline{toc}{subsection}{11.C.7 Tooling and Environment Supports for Recovery Execution}

Category-linked posture (no shopping lists; capability-level only; timing must not precede Stage 6):

\begin{itemize}
\item 7.11 Footwear Systems and Foot Care:
\begin{itemize}
\item stable footwear/sock interface; blister and hotspot posture; drying and rotation discipline.
\end{itemize}
\item 7.12 Recovery, Mobility, and Tissue-Care Implements:
\begin{itemize}
\item mat and basic mobility/tissue-care tools with safe-use limits; no “fix injury” claims.
\end{itemize}
\item 7.06 Recovery, Sleep, and Stress-Support Aids (optional tools):
\begin{itemize}
\item OTC-only posture; behavior-first hierarchy; no novelty near gates; no dosing in this deliverable.
\end{itemize}
\item 7.20 Personal Hygiene, Health, and Sanitation Items:
\begin{itemize}
\item hygiene and skin integrity consumables; basic coverage supplies; nail care.
\end{itemize}
\item 7.19 Protective and Assistive Equipment:
\begin{itemize}
\item visibility controls for low-light/traffic-adjacent sessions; prevention posture; sleeves comfort-only.
\end{itemize}
\item 7.10 Logistics, Maintenance, Storage, and Sustainment:
\begin{itemize}
\item drying workflows, staging, and inspection cadence to reduce preventable failures.
\end{itemize}
\end{itemize}

Rules (binding):

\begin{itemize}
\item capability timing \textbf{MUST} NOT be asserted earlier than Stage 6 introduces the category at Tier 1+.
\item missing tools do not justify unsafe substitutes; default to safer exposures and reslot gate events.
\item if any timing mismatch is detected between needed tooling and Stage 6 timing, flag "Requires Stage 8 review" here.
\end{itemize}

\begingroup
\scriptsize
\setlength{\tabcolsep}{2pt}
\renewcommand{\arraystretch}{1.12}

\begin{longtblr}{
width=\linewidth,
colspec = {
X[c,1.70]
X[c,1.20]
X[c,3.05]
X[c,3.05]
},
rowhead=1,
hlines,
vlines,
cells = {hsep=2pt, vsep=6pt, halign=c, valign=m},
row{1} = {bg=StageBlue!15, fg=black, font=\bfseries, valign=m},
row{odd} = {bg=white},
row{even} = {bg=LightGray},
}
\Hdr{Support Area} &
\Hdr{Category Reference} &
\Hdr{Minimum Operational Capability} &
\Hdr{If Missing Action} \\
Foot integrity &
7.11; 7.20 &
stable footwear/sock system; hotspot monitoring; drying posture &
downshift \textbf{RUN}/\textbf{RUCK} exposure; reslot tests; \textbf{HOLD} if gate evidence depends on foot integrity \\
Mobility/recovery implements &
7.12 &
mat + basic mobility/tissue-care tools; safe-use limits &
substitute with household alternatives where safe; preserve cadence; log limitations \\
Sleep support posture &
7.06 (optional) &
behavior-first sleep hygiene posture; no novelty &
do not add new aids near gates; rely on routine stability; consider maintenance reset per Stage 10 if collapse \\
Hygiene/sanitation &
7.20 &
clean-to-ready posture; skin integrity consumables; nail care &
increase hygiene cadence; avoid wet-gear loops; downshift if skin breakdown occurs \\
Visibility/protection &
7.19 &
session-level visibility controls when conditions require; hand protection as prevention &
substitute to daylight/indoor; do not proceed in low-light/traffic-adjacent conditions \\
Staging/drying/logistics &
7.10 &
drying workflow; staged kit readiness; inspection cadence &
treat as sustainment gap; increase staging discipline; reslot tests if readiness cannot be assured \\
\end{longtblr}

\endgroup

Requires Stage 8 review flags (if encountered):

\begin{itemize}
\item Any phase demand that requires a tooling capability earlier than Stage 6 tier timing provides.
\item Any attempt to treat optional tools (7.06) as required for deployability.
\end{itemize}

\newpage

\subsection*{11.C.8 Compliance Checklist — Pass/Fail}
\phantomsection
\addcontentsline{toc}{subsection}{11.C.8 Compliance Checklist — Pass/Fail}

\begin{itemize}
\item \textbf{PASS}/\textbf{FAIL}: No programming detail included (no workouts, sets/reps, intervals, weekly splits, or day-by-day schedules).
\item \textbf{PASS}/\textbf{FAIL}: No new tests, thresholds, domains, phase purposes, or equipment categories introduced.
\item \textbf{PASS}/\textbf{FAIL}: \textbf{STOP} \textbf{NOW} and \textbf{MEDICAL} \textbf{HOLD} posture stated and aligned to Stage 0.5 / Stage 4.F governance; Stage 1.F tokens used exactly.
\item \textbf{PASS}/\textbf{FAIL}: Downshift and reslot posture aligns to Stage 3.E; no stacking make-ups is explicit; \textbf{INVALID} tests provide no gate credit.
\item \textbf{PASS}/\textbf{FAIL}: Documentation requirements align to Stage 2.E missing-data doctrine (\textbf{UNKNOWN}/\textbf{MISSING}; no inference) and decision-grade records posture.
\item \textbf{PASS}/\textbf{FAIL}: Stage 7 category references are used correctly without shopping guidance; capability timing is not asserted earlier than Stage 6.
\item \textbf{PASS}/\textbf{FAIL}: Any mismatch is flagged "Requires Stage 8 review" in 11.C.6 and 11.C.7 rather than silently fixed.
\end{itemize}

\end{document}